\section{Mais resultados}
\def\sympathA{optimal-import-petro-layer-260514-14h05m22s}
\def\sympathB{optimal-import-petro-layer-blob-260514-14h05m28s}
\def\sympathC{optimal-import-petro-layer-layer-diag-260514-14h05m36s}
\def\sympathD{optimal-import-petro-layer-layer-column-260514-14h05m32s}
\def\sympathE{optimal-import-petro-layer-bracket-260514-14h05m43s}
\def\sympathF{optimal-import-petro-surround-260514-14h05m24s}
\def\sympathG{optimal-import-petro-layer-scolumn-layer-260514-14h05m40s}
\def\sympathH{optimal-import-petro-layer-layer-scolumn-260514-14h05m38s}
\def\sympathI{optimal-import-petro-semi-surround-260514-14h05m42s}
\def\sympathJ{optimal-import-petro-layer-cap-260514-14h05m45s}

\definecolor{ColA}{HTML}{E0BB00}
\definecolor{ColB}{HTML}{86AF00}
\definecolor{ColC}{HTML}{13E000}
\definecolor{ColD}{HTML}{00C133}
\definecolor{ColE}{HTML}{00CCCC}
\definecolor{ColF}{HTML}{0066FF}
\definecolor{ColG}{HTML}{3300FF}
\definecolor{ColH}{HTML}{CC00FF}
\definecolor{ColI}{HTML}{FF0099}
\definecolor{ColJ}{HTML}{FF0000}
\def\|{\text{\textbar}}
\tikzstyle{resA} = [
thick, ColA, dash pattern = {on 1pt off 0pt},
%mark=square*, mark size=2, mark options={
%	rotate=0,
%	solid,
%	opacity=0.5,
%	fill opacity=0.5
%}
]
\tikzstyle{resB} = [
thick, ColB, dash pattern = {on 1pt off 1pt},
%mark=triangle*, mark size=3, mark options={
%	rotate=0,
%	solid,
%	opacity=0.5,
%	fill opacity=0.5
%}
]
\tikzstyle{resC} = [
thick, ColC, dash pattern = {on 4pt off 2pt},
%mark=star, mark size=3, mark options={
%	rotate=0,
%	solid,
%	very thick,
%	opacity=0.5,
%	fill opacity=0.5
%}
]
\tikzstyle{resD} = [
thick, ColD, dash pattern = {on 4pt off 2pt on 1pt off 2pt},
%mark=square*, mark size=2, mark options={
%	rotate=45,
%	solid,
%	opacity=0.5,
%	fill opacity=0.5
%}
]
\tikzstyle{resE} = [
thick, ColE, dash pattern = {on 3pt off 1pt},
%mark=triangle*, mark size=1.5, mark options={
%	rotate=180,
%	solid,
%	fill opacity=0.5
%}
]
\tikzstyle{resF} = [
thick, ColF, dash pattern = {on 3pt off 2pt},
%mark=x, mark size=2, mark options={
%	rotate=0,
%	solid,
%	very thick,
%	fill opacity=0.5
%}
]
\tikzstyle{resG} = [
thick, ColG, dash pattern = {on 3pt off 4pt},
%mark=+, mark size=2, mark options={
%	rotate=0,
%	solid,
%	very thick,
%	fill opacity=0.5
%}
]
\tikzstyle{resH} = [
thick, ColH, dash pattern = {on 3pt off 4pt},
%mark=diamond*, mark size=2, mark options={
%	rotate=90,
%	solid,
%	very thick,
%	fill opacity=0.5
%}
]
\tikzstyle{resI} = [
thick, ColI, dash pattern = {on 3pt off 4pt},
%mark=pentagon*, mark size=2, mark options={
%	rotate=0,
%	solid,
%	very thick,
%	fill opacity=0.5
%}
]
\tikzstyle{resJ} = [
thick, ColJ, dash pattern = {on 3pt off 4pt},
%mark=*, mark size=2, mark options={
%	rotate=0,
%	solid,
%	very thick,
%	fill opacity=0.5
%}
]
\tikzstyle{resTTT} = [
	mark=square*, mark size=2, mark options={
		rotate=45,
		solid,
		opacity=0.5,
		fill opacity=0.5
	}
]
\tikzstyle{resSVD} = [
	mark=triangle*, mark size=3, mark options={
		rotate=0,
		solid,
		opacity=0.5,
		fill opacity=0.5
	}	
]

Como sugerido, os métodos utilizados na \cref{sec:simulation_and_results} (TTT e SVD) serão reaplicados aqui, em uma gama maior de condições de dispositivos. Serão testadas as seguintes configurações:

\begin{itemize}
	\item Camada\|Camada (C\|C) - Idêntico à configuração em camadas da \cref{sec:simulation_and_results}.
	\item Camada\|Bolha (C\|B) - Fontes no leito em uma camada, e receptores em 3 camadas próximas à superfície de 9 dispositivos cada, emulando várias viagens de um ROV.
	\item Camada\|Camada+Diagonal (C\|CD) - Fontes no leito em uma camada, e duas viagens de receptores: uma camada horizontal próxima à superfície, e uma viagem na diagonal até $500\si{\meter}$.
	\item Camada\|Camada+Apêndice (C\|CA) - Fontes no leito em uma camada, e duas viagens de receptores: uma camada horizontal próxima à superfície, e uma vertical até $500\si{\meter}$.
	\item Camada\|Garra (C\|G) - Fontes no leito em uma camada, receptores em uma cada próxima à superfície, e receptores em dois apêndices próximos às laterais da viagem horizontal.
	
	\item Camada+Coluna\|Camada+Coluna (CC\|CC) - Semelhante à configuração no entorno da \cref{sec:simulation_and_results}, mas com mais dispositivos.
	\item Camada+Vertical\|Camada (CV\|C) - Fontes no leito em uma camada + alguns (poucos) verticalmente, e receptores na superfície em uma camada.
	\item Camada\|Camada+Vertical (C\|CV) - Fontes no leito em uma camada, e receptores na superfície em uma camada + alguns (poucos) verticalmente.
	\item Camada+Vertical\|Camada+Vertical (CV\|CV) - Fontes no leito em uma camada + alguns (poucos) verticalmente, e receptores na superfície em uma camada + alguns (poucos) verticalmente.
	\item Camada\|Camada+Verticais (C\|CVs) - Fontes no leito em uma camada, e receptores na superfície em uma camada + duas colunas.
\end{itemize}

As configurações referentes a cada um desses casos encontram-se na \cref{fig:sec7:dispositivos}. Note que as configurações apresentam quantidades diferentes de dispositivos, em especial de fontes. Apesar de ser mais ``justo'' comparar configurações com quantidades iguais de dispositivos, a escolha aqui feita foi pensando na prática, em que não seria necessário reduzir a quantidade de medições (receptores) horizontais para medir na vertical, por exemplo.

\input{figures/A.a.ssf_discretizado}

Além disso, para facilitar visualização, os resultados serão apresentados para dois grupos diferentes:

\paragraph{Grupo 1)} consiste das 5 primeiras configurações (C\|C, C\|B, C\|CD, C\|CA, C\|G), em que os receptores permanecem exclusivamente próximo à superfície (acima de $500\si{\meter}$), e as fontes estão exclusivamente no leito.

\paragraph{Grupo 2)} consiste das 5 últimas configurações (CC\|CC, CV\|C, C\|CV, CV\|CV, C\|CVs), em que ambos os dispositivos podem ser alocados ao longo de toda a profundidade.

\begin{figure}
    \centering\setlength{\parindent}{0pt}
    \begin{subfigure}[t]{0.45\linewidth}
        \centering
            \begin{tikzpicture}
                \begin{lineplot}{Iterações}{$\rmse_{t}$ ($\si{\milli\second}$)}[legend to name=lineplots3, xmin=0.5, xmax=10.5, ymode=log,
                	legend style={
                		fill=white,
                		%			fill opacity=0.5.6,
                		draw opacity=1,
                		text opacity=1,
                		%			at={(1.05, 0.5)},
                		%			anchor=west,
                		legend columns=4,
                		/tikz/every even column/.append style={column sep=1em}
                	},]
                    {
                    	\def\sympath{\sympathA}
	                    \addplot[draw, resA, resSVD] table[x=x,y=y, col sep=comma] {figures/data/\sympath/metric_time_svd.csv};  
	                    \addplot[draw, resA, resTTT] table[x=x,y=y, col sep=comma] {figures/data/\sympath/metric_time_classic.csv}; 
                    }
                    {
                    	\def\sympath{\sympathB}
                    	\addplot[draw, resB, resSVD] table[x=x,y=y, col sep=comma] {figures/data/\sympath/metric_time_svd.csv};  
                    	\addplot[draw, resB, resTTT] table[x=x,y=y, col sep=comma] {figures/data/\sympath/metric_time_classic.csv}; 
                    }
                    
                    {
                    	\def\sympath{\sympathC}
                    	\addplot[draw, resC, resSVD] table[x=x,y=y, col sep=comma] {figures/data/\sympath/metric_time_svd.csv};  
                    	\addplot[draw, resC, resTTT] table[x=x,y=y, col sep=comma] {figures/data/\sympath/metric_time_classic.csv}; 
                    }
                    
                    {
                    	\def\sympath{\sympathD}
                    	\addplot[draw, resD, resSVD] table[x=x,y=y, col sep=comma] {figures/data/\sympath/metric_time_svd.csv};  
                    	\addplot[draw, resD, resTTT] table[x=x,y=y, col sep=comma] {figures/data/\sympath/metric_time_classic.csv}; 
                    }
                    
                    {
                    	\def\sympath{\sympathE}
                    	\addplot[draw, resE, resSVD] table[x=x,y=y, col sep=comma] {figures/data/\sympath/metric_time_svd.csv};  
                    	\addplot[draw, resE, resTTT] table[x=x,y=y, col sep=comma] {figures/data/\sympath/metric_time_classic.csv}; 
                    }
                    
                    {
                    	\def\sympath{\sympathF}
                    	\addplot[draw, resF, resSVD] table[x=x,y=y, col sep=comma] {figures/data/\sympath/metric_time_svd.csv};  
                    	\addplot[draw, resF, resTTT] table[x=x,y=y, col sep=comma] {figures/data/\sympath/metric_time_classic.csv}; 
                    }
                    {
                    	\def\sympath{\sympathG}
                    	\addplot[draw, resG, resSVD] table[x=x,y=y, col sep=comma] {figures/data/\sympath/metric_time_svd.csv};  
                    	\addplot[draw, resG, resTTT] table[x=x,y=y, col sep=comma] {figures/data/\sympath/metric_time_classic.csv}; 
                    }
                    
                    {
                    	\def\sympath{\sympathH}
                    	\addplot[draw, resH, resSVD] table[x=x,y=y, col sep=comma] {figures/data/\sympath/metric_time_svd.csv};  
                    	\addplot[draw, resH, resTTT] table[x=x,y=y, col sep=comma] {figures/data/\sympath/metric_time_classic.csv}; 
                    }
                    
                    {
                    	\def\sympath{\sympathI}
                    	\addplot[draw, resI, resSVD] table[x=x,y=y, col sep=comma] {figures/data/\sympath/metric_time_svd.csv};  
                    	\addplot[draw, resI, resTTT] table[x=x,y=y, col sep=comma] {figures/data/\sympath/metric_time_classic.csv}; 
                    }
                    
                    {
                    	\def\sympath{\sympathJ}
                    	\addplot[draw, resJ, resSVD] table[x=x,y=y, col sep=comma] {figures/data/\sympath/metric_time_svd.csv};  
                    	\addplot[draw, resJ, resTTT] table[x=x,y=y, col sep=comma] {figures/data/\sympath/metric_time_classic.csv}; 
                    }

                    % \addplot[draw, resA]
                    \addlegendentry{SVD - C\|C};
                    \addlegendentry{TTT - C\|C};
                    \addlegendentry{SVD - C\|B};
                    \addlegendentry{TTT - C\|B};
                    \addlegendentry{SVD - C\|CD};
                    \addlegendentry{TTT - C\|CD};
                    \addlegendentry{SVD - C\|CA};
                    \addlegendentry{TTT - C\|CA};
                    \addlegendentry{SVD - C\|G};
                    \addlegendentry{TTT - C\|G};
                    %
                    \addlegendentry{SVD - CC\|CC};
                    \addlegendentry{TTT - CC\|CC};
                    \addlegendentry{SVD - CV\|C};
                    \addlegendentry{TTT - CV\|C};
                    \addlegendentry{SVD - C\|CV};
                    \addlegendentry{TTT - C\|CV};
                    \addlegendentry{SVD - CV\|CV};
                    \addlegendentry{TTT - CV\|CV};
                    \addlegendentry{SVD - C\|CVs};
                    \addlegendentry{TTT - C\|CVs};
                \end{lineplot}
            \end{tikzpicture}
        \caption{Erro RMS de tempo de trânsito.}
        \label{fig:sec7:tt_ssf_rms_errors:subfig1}
    \end{subfigure} \hspace*{2em}
    \begin{subfigure}[t]{0.45\linewidth}
        \centering
            \begin{tikzpicture}
                \begin{lineplot}{Iterações}{$\rmse_{\ssf}$ ($\si{\meter/\second}$)}[xmin=0.5, xmax=10.5, ymode=log,]
                    {
                    	\def\sympath{\sympathA}
                    	\addplot[draw, resA, resSVD] table[x=x,y=y, col sep=comma] {figures/data/\sympath/metric_ssf_svd.csv};  
                    	\addplot[draw, resA, resTTT] table[x=x,y=y, col sep=comma] {figures/data/\sympath/metric_ssf_classic.csv}; 
                    }
                    {
                    	\def\sympath{\sympathB}
                    	\addplot[draw, resB, resSVD] table[x=x,y=y, col sep=comma] {figures/data/\sympath/metric_ssf_svd.csv};  
                    	\addplot[draw, resB, resTTT] table[x=x,y=y, col sep=comma] {figures/data/\sympath/metric_ssf_classic.csv}; 
                    }
                    
                    {
                    	\def\sympath{\sympathC}
                    	\addplot[draw, resC, resSVD] table[x=x,y=y, col sep=comma] {figures/data/\sympath/metric_ssf_svd.csv};  
                    	\addplot[draw, resC, resTTT] table[x=x,y=y, col sep=comma] {figures/data/\sympath/metric_ssf_classic.csv}; 
                    }
                    
                    {
                    	\def\sympath{\sympathD}
                    	\addplot[draw, resD, resSVD] table[x=x,y=y, col sep=comma] {figures/data/\sympath/metric_ssf_svd.csv};  
                    	\addplot[draw, resD, resTTT] table[x=x,y=y, col sep=comma] {figures/data/\sympath/metric_ssf_classic.csv}; 
                    }
                    
                    {
                    	\def\sympath{\sympathE}
                    	\addplot[draw, resE, resSVD] table[x=x,y=y, col sep=comma] {figures/data/\sympath/metric_ssf_svd.csv};  
                    	\addplot[draw, resE, resTTT] table[x=x,y=y, col sep=comma] {figures/data/\sympath/metric_ssf_classic.csv}; 
                    }
                    
                    {
                    	\def\sympath{\sympathF}
                    	\addplot[draw, resF, resSVD] table[x=x,y=y, col sep=comma] {figures/data/\sympath/metric_ssf_svd.csv};  
                    	\addplot[draw, resF, resTTT] table[x=x,y=y, col sep=comma] {figures/data/\sympath/metric_ssf_classic.csv}; 
                    }
                    {
                    	\def\sympath{\sympathG}
                    	\addplot[draw, resG, resSVD] table[x=x,y=y, col sep=comma] {figures/data/\sympath/metric_ssf_svd.csv};  
                    	\addplot[draw, resG, resTTT] table[x=x,y=y, col sep=comma] {figures/data/\sympath/metric_ssf_classic.csv}; 
                    }
                    
                    {
                    	\def\sympath{\sympathH}
                    	\addplot[draw, resH, resSVD] table[x=x,y=y, col sep=comma] {figures/data/\sympath/metric_ssf_svd.csv};  
                    	\addplot[draw, resH, resTTT] table[x=x,y=y, col sep=comma] {figures/data/\sympath/metric_ssf_classic.csv}; 
                    }
                    
                    {
                    	\def\sympath{\sympathI}
                    	\addplot[draw, resI, resSVD] table[x=x,y=y, col sep=comma] {figures/data/\sympath/metric_ssf_svd.csv};  
                    	\addplot[draw, resI, resTTT] table[x=x,y=y, col sep=comma] {figures/data/\sympath/metric_ssf_classic.csv}; 
                    }
                    
                    {
                    	\def\sympath{\sympathJ}
                    	\addplot[draw, resJ, resSVD] table[x=x,y=y, col sep=comma] {figures/data/\sympath/metric_ssf_svd.csv};  
                    	\addplot[draw, resJ, resTTT] table[x=x,y=y, col sep=comma] {figures/data/\sympath/metric_ssf_classic.csv}; 
                    }
                \end{lineplot}
            \end{tikzpicture}
        \caption{Erro RMS do SSF.}
        \label{fig:sec7:tt_ssf_rms_errors:subfig2}
    \end{subfigure}

    \vspace*{.5em}
    \colorbar{lineplots3}
    \caption{Evolução dos erros RMS do processo de estimação do SSF, para os métodos TTT e SVD, nos cenários de dispositivos no perímetro e em camadas, ao longo de 10 iterações.}
    \label{fig:sec7:tt_ssf_rms_errors}
\end{figure}